\begin{beginnerstatement}[Kurz erklärt: Wo Nutzerdaten gespeichert werden]

Programmdateien bestimmen, was die Anwendung kann. Einstellungen wie eine
Serveradresse oder ein Datenbankzugang bestimmen, wie sie gestartet wird.
Nutzerdaten entstehen dagegen beim Arbeiten mit dem Produkt, etwa Dokumente,
Planungsstände oder Bilder. Nicht jede App braucht dafür den gleichen
Speicher. Kleine lokale Anwendungen können Dateien verwenden; umfangreichere
lokale Anwendungen können eine eingebettete Datenbank benötigen.
Serverbasierte Anwendungen können eine zentrale Datenbank nutzen, während
Bilder und Dokumente separat gespeichert werden können. Das Template wählt
diese Lösung bewusst nicht automatisch. Das Produktprojekt muss entscheiden,
welcher Speicher maßgeblich ist, und bei mehreren Speichern zusätzlich Regeln
für Offline-Betrieb, Synchronisation und Konflikte festlegen. Die Baseline
enthält noch keine lokale Speicher- oder Synchronisationslösung.

\end{beginnerstatement}
